% Vietnamese translation for OI Public Library
% Source-File: IOI中国国家候选队论文/2008/Day2/5.陈瑜希《Pólya计数法的应用》/pálya计数法的应用.ppt
\documentclass[11pt]{article}
\usepackage{oipl-vn}

\title{Bản trình chiếu: Ứng dụng của phương pháp đếm Pólya}
\author{Chen Yuxi}
\date{2008}

\begin{document}
\maketitle

\origtitle{Slide companion for applications of Pólya counting}
\sourcepath{IOI China candidate team papers / 2008 / Day 2 / Chen Yuxi.ppt}

\section{Phạm vi bản dịch}

Tệp PowerPoint đi cùng bài viết đã được dịch từ PDF. Bản ghi chú này giữ lại
các khái niệm cần thiết để dùng Burnside và Pólya trong bài toán đếm có đối
xứng.

\section{Vấn đề đối xứng}

Khi các cấu hình được xem là giống nhau sau phép quay, lật hoặc hoán vị, cách
đếm trực tiếp thường đếm trùng. Ví dụ vòng gồm các ký tự có nhiều biểu diễn
tuyến tính khác nhau nhưng cùng một cấu hình trên vòng tròn.

\section{Burnside và Pólya}

Định lý Burnside nói rằng số quỹ đạo bằng trung bình số cấu hình bất động qua
các phép đối xứng của nhóm. Phương pháp Pólya mở rộng ý tưởng này bằng cách
dùng chu trình của hoán vị để tính số cách tô màu hoặc gán nhãn thỏa đối xứng.

\section{Kết luận}

Trong thi lập trình, phương pháp Pólya thường xuất hiện khi đề có các từ khóa
như vòng, đa giác, xoay, lật, tô màu hoặc cấu hình giống nhau về bản chất.
Việc khó nhất là xác định đúng nhóm đối xứng và số chu trình của từng phép
biến đổi.

\end{document}
